\documentclass[11pt]{letter}
\usepackage[T1]{fontenc}
\usepackage{lmodern}
\usepackage[margin=1in]{geometry}

\address{School of Electronic and Information Engineering\\
Suzhou University of Science and Technology\\
Suzhou 215009, Jiangsu, China}

\begin{document}

\begin{letter}{The Editor-in-Chief\\International Journal of Refrigeration}

\opening{Dear Editor,}

We are pleased to submit our manuscript entitled ``Event-Triggered Predictive Reinforcement Learning with Unsupervised Dynamic Event Gating for Energy Efficient HVAC Control'' for consideration as a Research Article in the \textit{International Journal of Refrigeration}.

This study addresses energy-efficient operation of the chilled-water side of a commercial-building cooling system, where the chiller is the dominant energy-consuming component. Fixed-interval reinforcement learning controllers execute policies uniformly even though building cooling loads evolve at different rates, causing redundant setpoint actuation during stable operation and delayed response to load changes. We propose event-triggered predictive reinforcement learning with unsupervised dynamic event gating (ET-PRL), which converts fixed-step policy execution into state-dependent action updates by treating control triggering as online anomaly detection over a predictive state. A multi-scale fused anomaly score and a local--global two-layer adaptive threshold keep the trigger boundary sensitive to abrupt disturbances while adapting to intraday variation and long-term operating-condition drift. The method is evaluated in a calibrated EnergyPlus model of a commercial building driven by measured chiller operation data. Compared with a fixed-step reinforcement learning controller using the same policy network, ET-PRL reduces policy executions by 53.54\% and daily chiller energy use by 2.13\%.

We believe this work fits the \textit{International Journal of Refrigeration} closely. The journal's aims and scope include refrigeration and air-conditioning systems, refrigerant and chilled-water systems, heat pumps, and the energy performance of building cooling systems. Our study directly targets the chilled-water supply-temperature control of a chiller plant, i.e., a core operating variable of the refrigeration cycle that governs chiller part-load efficiency and the overall cooling-system energy consumption. The reported daily chiller energy reduction, together with the reduced setpoint actuation that lowers compressor and pump cycling and actuator wear, is a practical contribution to more energy-efficient refrigeration and HVAC system operation.

The manuscript has six authors, all of whom have read and approved the final version and agree with its submission to the \textit{International Journal of Refrigeration}. The CRediT authorship contribution statement is included in the manuscript.

This manuscript is original, has not been published previously, and is not under consideration by any other journal. The authors declare that they have no competing interests. Generative AI-assisted tools were used for language editing and formatting support only; the authors reviewed and edited the content and take full responsibility for the publication.

Thank you for considering our submission.

\vspace{2\baselineskip}
\noindent Sincerely,\\[2\baselineskip]
Qiming Fu\\
School of Electronic and Information Engineering\\
Suzhou University of Science and Technology\\
Suzhou 215009, Jiangsu, China\\
E-mail: fqm@mail.usts.edu.cn\\[1\baselineskip]
Jianping Chen\\
School of Electronic and Information Engineering\\
Suzhou University of Science and Technology\\
Suzhou 215009, Jiangsu, China\\
E-mail: alan@mail.usts.edu.cn

\end{letter}
\end{document}
